% Шаблон презентации доклада (Beamer)
% Лаборатория математики ФН1
% Компиляция: pdflatex presentation.tex && pdflatex presentation.tex

\documentclass[aspectratio=169, 12pt]{beamer}

\usepackage[T2A]{fontenc}
\usepackage[utf8]{inputenc}
\usepackage[russian]{babel}
\usepackage{amsmath, amssymb}
\usepackage{graphicx}
\usepackage{booktabs}

\usetheme{Madrid}
\usecolortheme{whale}
\setbeamertemplate{navigation symbols}{}
\setbeamertemplate{caption}[numbered]

% Фирменные цвета лаборатории
\definecolor{labnavy}{RGB}{18, 58, 95}
\definecolor{labaccent}{RGB}{255, 200, 87}
\setbeamercolor{structure}{fg=labnavy}
\setbeamercolor{alerted text}{fg=labnavy}

\title[Краткое название]{Название доклада}
\author[Фамилия И.\,О.]{Фамилия Имя Отчество}
\institute[ФН1]{Лаборатория математики ФН1}
\date{\today}

\begin{document}

\begin{frame}[plain]
  \titlepage
\end{frame}

% ---------- Мотивация ----------
\begin{frame}{Зачем эта задача}
  \begin{itemize}
    \item где встречается на практике;
    \item почему известные подходы не дают ответа;
    \item что хотелось бы получить.
  \end{itemize}
\end{frame}

% ---------- Постановка ----------
\begin{frame}{Постановка задачи}
  Рассматривается
  \begin{equation}\label{eq:problem}
    L u = f, \qquad u \in D(L),
  \end{equation}
  где $L$~--- линейный оператор в гильбертовом пространстве $H$.

  \vspace{0.6em}
  \alert{Требуется:} найти условия однозначной разрешимости~\eqref{eq:problem}
  и оценить норму решения.
\end{frame}

% ---------- Известные результаты ----------
\begin{frame}{Что известно}
  \begin{itemize}
    \item[{[1]}] Автор А.\,Б. (2019) --- результат при условии $\ldots$
    \item[{[2]}] Автор В.\,Г. (2023) --- обобщение на случай $\ldots$
  \end{itemize}

  \vspace{0.6em}
  Открытым остаётся случай $\ldots$
\end{frame}

% ---------- Основной результат ----------
\begin{frame}{Основной результат}
  \begin{block}{Теорема}
    Формулировка. Условия перечисляются явно, обозначения --- те же,
    что на предыдущем слайде.
  \end{block}

  \vspace{0.6em}
  \begin{itemize}
    \item чем результат отличается от известных;
    \item какова цена ослабления условий.
  \end{itemize}
\end{frame}

% ---------- Эксперимент ----------
\begin{frame}{Численный эксперимент}
  \begin{columns}[T]
    \begin{column}{0.5\textwidth}
      % \includegraphics[width=\textwidth]{plot.pdf}
      \centering\footnotesize (график сходимости)
    \end{column}
    \begin{column}{0.5\textwidth}
      \centering
      \begin{tabular}{ccc}
        \toprule
        $h$ & Ошибка & Порядок \\
        \midrule
        $0{,}1$  & $2{,}1\cdot10^{-3}$ & --- \\
        $0{,}05$ & $5{,}3\cdot10^{-4}$ & $1{,}99$ \\
        $0{,}025$& $1{,}3\cdot10^{-4}$ & $2{,}03$ \\
        \bottomrule
      \end{tabular}
    \end{column}
  \end{columns}

  \vspace{0.8em}
  Наблюдаемый порядок согласуется с теоретической оценкой $O(h^2)$.
\end{frame}

% ---------- Выводы ----------
\begin{frame}{Выводы}
  \begin{enumerate}
    \item получено $\ldots$;
    \item показано, что $\ldots$;
    \item в дальнейшем планируется $\ldots$
  \end{enumerate}

  \vspace{1em}
  \centering\alert{Спасибо за внимание!}
\end{frame}

% ---------- Литература ----------
\begin{frame}[allowframebreaks]{Литература}
  \footnotesize
  \begin{thebibliography}{9}
    \bibitem{ref1} Автор А.\,Б. Название статьи~// Журнал.~--- 2019.~--- Т.~5, №~1.~--- С.~3--17.
    \bibitem{ref2} Автор В.\,Г. Название книги.~--- М.: Издательство, 2023.
  \end{thebibliography}
\end{frame}

\end{document}
